\chapter{Einleitung}\label{ch:einleitung}

Die Einleitung führt in das Thema der Arbeit ein und beschreibt den fachlichen sowie praktischen Kontext.
Sie soll die Leserinnen und Leser abholen, die Problemstellung herleiten und das Ziel klar formulieren.

Ein typischer Aufbau umfasst:

\begin{enumerate}
\item \textbf{Motivation und Kontext:} Warum ist das Thema wichtig? Wo ist die Anwendung relevant? Beispielhafte Problemfelder oder Anwendungsgebiete.
\item \textbf{Problemstellung:} Was ist bisher ungelöst oder verbesserungswürdig? Formulieren Sie eine konkrete Forschungsfrage oder ein Entwicklungsziel.
\item \textbf{Zielsetzung und Beitrag:} Welche konkrete Aufgabe bearbeitet die Arbeit, und welchen Mehrwert liefert sie? Wird ein Modell entwickelt, ein Algorithmus evaluiert, ein Prototyp konstruiert?
\item \textbf{Aufbau der Arbeit:} Kurze Beschreibung, was in den einzelnen Kapiteln behandelt wird.
\end{enumerate}

\textit{Beispiel:} „Kapitel 2 stellt die theoretischen Grundlagen vor, Kapitel 3 beschreibt die Methodik, Kapitel 4 zeigt die Ergebnisse und Kapitel 5 diskutiert diese.
Kapitel 6 fasst die Arbeit zusammen und gibt einen Ausblick.“



